% =============================================================================
% Write Path in OLTP Workloads — rootconf 2026
% A hands-on PostgreSQL workshop
%
% Compile with:  lualatex slides.tex   (run twice for page numbers)
% =============================================================================
\documentclass[aspectratio=169,10pt]{beamer}

% ── Theme ─────────────────────────────────────────────────────────────────────
\usetheme[
  progressbar=frametitle,
  block=fill,
  sectionpage=progressbar,
  numbering=fraction,
]{metropolis}
\usepackage{appendixnumberbeamer}

% ── Fonts (LuaLaTeX) ──────────────────────────────────────────────────────────
% Metropolis uses Fira Sans / Fira Mono, both included in MacTeX.
\usepackage{fontspec}

% ── PostgreSQL colour palette ─────────────────────────────────────────────────
\definecolor{pgblue}{HTML}{336791}    % PostgreSQL brand blue
\definecolor{pgdark}{HTML}{1e3a5c}
\definecolor{pglight}{HTML}{eaf1f8}
\definecolor{pgaccent}{HTML}{0075b8}
\definecolor{codebg}{HTML}{f4f4f4}
\definecolor{alertred}{HTML}{c0392b}

\setbeamercolor{frametitle}{bg=pgblue, fg=white}
\setbeamercolor{progress bar}{fg=pgaccent, bg=pgblue!30}
\setbeamercolor{alerted text}{fg=pgaccent}
\setbeamercolor{block title}{bg=pgblue!15, fg=pgdark}
\setbeamercolor{block title alerted}{bg=alertred!15, fg=alertred!80!black}

% ── Packages ──────────────────────────────────────────────────────────────────
\usepackage{listings}
\usepackage{xcolor}
\usepackage{booktabs}
\usepackage{tabularx}
\usepackage{array}
\usepackage{fontawesome5}
\usepackage{fancyvrb}
\usepackage{hyperref}

% ── SQL listing style ─────────────────────────────────────────────────────────
\lstdefinestyle{sql}{
  language=SQL,
  basicstyle=\ttfamily\footnotesize,
  keywordstyle=\color{pgblue}\bfseries,
  commentstyle=\color{gray}\itshape,
  stringstyle=\color{teal},
  showstringspaces=false,
  breaklines=true,
  frame=single,
  backgroundcolor=\color{codebg},
  rulecolor=\color{pgblue!40},
  framesep=3pt,
  xleftmargin=2pt,
  xrightmargin=2pt,
  belowskip=0pt,
}

% ── ASCII / diagram listing style ─────────────────────────────────────────────
\lstdefinestyle{ascii}{
  basicstyle=\ttfamily\scriptsize,
  showstringspaces=false,
  frame=none,
  backgroundcolor=\color{codebg},
  belowskip=0pt,
  xleftmargin=2pt,
}

% ── Metadata ──────────────────────────────────────────────────────────────────
\title{Write Path in OLTP Workloads}
\subtitle{A Hands-On PostgreSQL Workshop}
\author{Amit Prabhudesai}
\institute{Tech Lead, Microsoft}
\date{rootconf 2026}

% =============================================================================
\begin{document}
% =============================================================================

\maketitle

% -----------------------------------------------------------------------------
% Slide: About Me
% -----------------------------------------------------------------------------
\begin{frame}{About Me}
  \begin{columns}[T]
    \column{0.62\textwidth}
    \begin{itemize}
      \item \textbf{Tech Lead, Microsoft}
      \item \href{https://github.com/documentdb/documentdb}{%
        \faGithub\enspace\texttt{documentdb/documentdb}}
      \item Ex-InMobi (China-engg all-up; Architect, InMobi Ad SDK)
      \item Worn both hats: Engineering Manager \textit{and} Tech Lead
    \end{itemize}
    \vspace{0.6em}
    {\small Slides, exercises, and devcontainer:}\\[0.3em]
    \href{https://github.com/amitprabhudesai/oltp-workloads-pg-workshop}{%
      \faGithub\enspace\texttt{amitprabhudesai/oltp-workloads-pg-workshop}}

    \column{0.34\textwidth}
    \vspace{1em}
    {\large\faGithub}\enspace\texttt{amitprabhudesai}\\[0.5em]
    {\large\faTwitter}\enspace\href{https://x.com/amitprabhudesai}{\texttt{@amitprabhudesai}}
  \end{columns}
\end{frame}

% -----------------------------------------------------------------------------
% Slide: About This Workshop
% -----------------------------------------------------------------------------
\begin{frame}{About This Workshop}

  \begin{block}{Database systems as fact-keepers}
    \textbf{Durability.} Records persist reliably ---
     committed data survives crashes and restarts.\\[0.3em]
    \textbf{Correctness.} Change is constant; interleaved reads and
     writes must faithfully reflect the intent of every update.\\[0.3em]
    \textbf{Performance.} Neither guarantee is useful with unbounded
     latency or unsustainable write amplification.\\[0.4em]
    {\small This workshop is about the mechanisms PostgreSQL uses
    to honour all three.}
  \end{block}
  \vspace{0.6em}
  \begin{columns}[T]
    \column{0.48\textwidth}
    \textbf{Module 1 --- WAL \& Durability}\\[0.2em]
    {\small How PostgreSQL guarantees durability without sacrificing throughput.}

    \column{0.48\textwidth}
    \textbf{Module 2 --- MVCC \& Concurrency}\\[0.2em]
    {\small How multiple writers coexist and what the hidden costs are.}
  \end{columns}

  \vspace{0.8em}
  \begin{center}
    {\small
    \faTools\enspace
    Each module: short theory \textrightarrow\ hands-on exercises \textrightarrow\ debrief\\[0.3em]
    \faDocker\enspace
    devcontainer included --- \texttt{docker} is enough, no local Postgres needed.
    }
  \end{center}
\end{frame}

% -----------------------------------------------------------------------------
% Slide: Why PostgreSQL?
% -----------------------------------------------------------------------------
\begin{frame}{Why PostgreSQL?}
  \begin{columns}[T]
    \column{0.58\textwidth}
    \begin{itemize}
      \item DocumentDB is a PostgreSQL extension
      \item \alert{\#1 most-used database} in the Stack\,Overflow Developer
            Survey --- \textbf{three years running} (2023, 2024, 2025)
      \item Well-documented internals: source is readable, WAL format is
            stable, behaviour is inspectable with plain SQL
      \item The concepts transfer directly:\\
            MySQL binlog \textbullet\ CockroachDB Raft log \textbullet\
            AWS Aurora redo log --- all WAL variants
      \item Open source: you can \textit{read what actually runs}
    \end{itemize}

    \column{0.38\textwidth}
    \vspace{0.4em}
    {\footnotesize Stack\,Overflow Developer Survey 2025\\
    \textit{Technology --- Databases}:}
    \vspace{0.3em}

    \begin{tabular}{lr}
      \toprule
      \textbf{Database} & \textbf{Usage} \\
      \midrule
      \textbf{PostgreSQL} & \textbf{49\,\%} \\
      MySQL               & 40\,\% \\
      SQLite              & 33\,\% \\
      Redis               & 28\,\% \\
      \bottomrule
    \end{tabular}

    \vspace{0.3em}
    {\scriptsize\textit{survey.stackoverflow.co/2025}}
  \end{columns}
\end{frame}

% -----------------------------------------------------------------------------
% Slide: Workshop Structure
% -----------------------------------------------------------------------------
\begin{frame}{Workshop Structure}
  \begin{columns}[T]
    \column{0.48\textwidth}
    \textbf{\faServer\enspace Module 1 --- WAL \& Durability}
    \begin{itemize}\small
      \item WAL write path \& record anatomy
      \item \texttt{pg\_stat\_wal}: observing the stream
      \item \texttt{synchronous\_commit}: the durability dial
      \item Checkpoints \& the recovery horizon
      \item Full-page images: the checkpoint tax
    \end{itemize}

    \column{0.48\textwidth}
    \textbf{\faDatabase\enspace Module 2 --- MVCC \& Concurrency}
    \begin{itemize}\small
      \item Multi-version heap pages
      \item Snapshot isolation in practice
      \item Lock contention \& row-level locking
      \item VACUUM \& bloat
    \end{itemize}
  \end{columns}

  \vspace{0.7em}
  \begin{center}
    {\small\faTerminal\enspace
     All exercises run in a self-contained devcontainer (PostgreSQL\,16 + psql).\\
     Open \texttt{modules/0X-.../exercises.sql} and follow along.}
  \end{center}
\end{frame}

% =============================================================================
\section{Module 1 --- WAL \& Durability}
% =============================================================================

% -----------------------------------------------------------------------------
% Slide: The Durability Problem — motivation
% -----------------------------------------------------------------------------
\begin{frame}{The Durability Problem}

  \begin{block}{}
    A \texttt{COMMIT} returns to the client.  The modified pages sit
    in shared buffers\footnote{\url{https://www.postgresql.org/docs/16/runtime-config-resource.html\#GUC-SHARED-BUFFERS}} --- in RAM.\\[0.4em]
    What happens if the server crashes before those pages reach disk?
  \end{block}

\end{frame}

% -----------------------------------------------------------------------------
% Slide: The Durability Problem — detail
% -----------------------------------------------------------------------------
\begin{frame}[fragile]{The Durability Problem}

  \begin{lstlisting}[style=ascii]
  (1) BEGIN; INSERT INTO tbl VALUES ('A'); COMMIT;
  (2)              BEGIN; INSERT INTO tbl VALUES ('B'); COMMIT; --> CRASH

  RAM:  +----------+        +----------+        +----------+
        | TABLE_A  |  --->  | TABLE_A  |  --->  | TABLE_A  |
        | [empty]  |        |   [ A ]  |        | [ B | A ]|
        +----------+        +----------+        +----------+
                                                      |
                                          never flushed to disk
  disk: +----------+
        | TABLE_A  |  <-- both commits LOST; page still empty after restart
        | [empty]  |
        +----------+
  \end{lstlisting}

  \vspace{0.4em}
  \textbf{The insight:} a \textit{log} is always appended sequentially.\footnote{Sequential writes are 10--100$\times$ cheaper than random writes on spinning disk; meaningfully cheaper even on SSD.}
  If the log is durable, the data pages can be written lazily.
\end{frame}

% -----------------------------------------------------------------------------
% Slide: WAL Principle
% -----------------------------------------------------------------------------
\begin{frame}{The Write-Ahead Log}

  \begin{block}{}
    Record all modifications as history data in persistent storage.
    This history is called \textbf{XLOG} or \textbf{WAL data}.
  \end{block}

\end{frame}

% -----------------------------------------------------------------------------
% Slide: Tuple Insertion with WAL
% -----------------------------------------------------------------------------
% picture first — let it speak
\begin{frame}[fragile,label=frTupleInsert]{Tuple Insertion with WAL}

\begin{lstlisting}[style=ascii]
  INSERT  --> (1) WAL buffer  : XLOG @ LSN_1  (write-ahead)
          --> (2) Shared Buffers: page modified, pd_lsn = LSN_1
  COMMIT  --> (3) pg_wal/     : write + fsync  (durable)

  INSERT  --> (1) WAL buffer  : XLOG @ LSN_2
          --> (2) Shared Buffers: page modified, pd_lsn = LSN_2
  COMMIT  --> (3) pg_wal/     : write + fsync

  -------------------------------------------------------

  Shared Buffers (RAM)          pg_wal/ (disk)
  +----------------------+       +----------------------+
  | TABLE_A              |       | XLOG @ LSN_2         |
  | [tup_1, tup_2]       |       | XLOG @ LSN_1         |
  | pd_lsn = LSN_2       |       | CKPT @ REDO_LSN  (0) |
  +----------------------+       +----------------------+
  lazy: flushed later             fsynced at COMMIT
  (bgwriter / checkpoint)         survives crash
\end{lstlisting}

{\tiny Heap tuple layout (header, \texttt{xmin}/\texttt{xmax}, data):
\hyperlink{frA3}{\beamergotobutton{Appendix A3}}\par}

\end{frame}

% text reinforces the picture
\begin{frame}{Tuple Insertion with WAL}

  \begin{enumerate}\small
    \setcounter{enumi}{-1}
    \item {\color{gray}\textbf{Checkpoint}: checkpointer writes a
          checkpoint record containing the latest \textbf{REDO point} to WAL.}
    \item \textbf{INSERT}: backend loads the page into shared buffers,
          inserts the tuple, writes an XLOG record at \texttt{LSN\_1},
          updates the page header: \texttt{pd\_lsn} \texttt{LSN\_0}
          $\to$ \texttt{LSN\_1}.
    \item \textbf{COMMIT}: WAL buffer flushed to \texttt{pg\_wal/}
          (\texttt{write} + \texttt{fsync}).
    \item Subsequent INSERT: XLOG record at \texttt{LSN\_2};
          \texttt{pd\_lsn} \texttt{LSN\_1} $\to$ \texttt{LSN\_2}.
    \item \textbf{CRASH}: shared buffers lost; \texttt{pg\_wal/}
          survives on disk. Recovery replays from the REDO point.
  \end{enumerate}

\end{frame}

% -----------------------------------------------------------------------------
% Slide: exec_simple_query() — WAL write path (also Appendix A6)
% -----------------------------------------------------------------------------
\begin{frame}[fragile]{\texttt{exec\_simple\_query()}: WAL Write Path}

\begin{lstlisting}[style=ascii]
exec_simple_query()                         [postgres.c]
  |
  +-- (1) ExtendCLOG()                      [clog.c]
  |        mark txid IN_PROGRESS in CLOG
  |
  +-- (2) heap_insert()                     [heapam.c]
  |        insert tuple into shared buffer page
  |   |
  |   +-- (3) XLogInsert()                  [xloginsert.c]
  |            write Heap/INSERT record to WAL buffer
  |            update page pd_lsn := LSN_1
  |
  +-- (4) finish_xact_command()             [postgres.c]
      |    commit action
      |
      +-- (4a) XLogInsert()                 [xloginsert.c]
      |         write Commit record to WAL buffer (LSN_2)
      |
      +-- (5) XLogWrite()                   [xlog.c]
      |        flush WAL buffer -> WAL segment file
      |        (fsync / fdatasync / O_DSYNC per wal_sync_method)
      |
      +-- (6) TransactionIdCommitTree()     [transam.c]
               mark txid COMMITTED in CLOG
\end{lstlisting}

\end{frame}

% -----------------------------------------------------------------------------
% Slide: Database Recovery Using WAL
% -----------------------------------------------------------------------------
\begin{frame}[fragile]{Database Recovery Using WAL}
\begin{lstlisting}[style=ascii]
  NORMAL OPERATION
  ---------------------------------------------------------
  Shared Buffers (RAM)       pg_wal/ (fsynced at COMMIT)
  +--------------------+     +-------------------------+
  | TABLE_A            |     | XLOG INSERT @ LSN_2     |
  | [tup_1, tup_2]     |     | XLOG INSERT @ LSN_1     |
  | pd_lsn = LSN_2     |     | CKPT        @ REDO_LSN  |
  +--------------------+     +-------------------------+

  ** CRASH **: server dies; shared buffers gone; pg_wal/ survives
  ---------------------------------------------------------
  +--------------------+     +-------------------------+
  | //// LOST /////    |     | XLOG INSERT @ LSN_2     |
  | (RAM, gone)        |     | XLOG INSERT @ LSN_1     |
  +--------------------+     | CKPT        @ REDO_LSN  |
                             +-------------------------+
  RECOVERY: replay pg_wal/ forward from REDO_LSN
  ---------------------------------------------------------
  LSN_1 > pd_lsn(0)?  YES --> apply insert;  pd_lsn = LSN_1
  LSN_2 > pd_lsn(1)?  YES --> apply insert;  pd_lsn = LSN_2
  LSN   <= pd_lsn?    NO  --> skip (already on disk)
  --> database restored to pre-crash state
\end{lstlisting}
\end{frame}

% -----------------------------------------------------------------------------
% Slide: Checkpoints motivation
% -----------------------------------------------------------------------------
\begin{frame}{Checkpoints}

  \begin{block}{}
    Recovery replays WAL from some starting point.
    If that starting point is the very beginning of WAL history,
    recovery time is unbounded.\\[0.4em]
    How does the database keep restart time predictable?
  \end{block}

\end{frame}

% -----------------------------------------------------------------------------
% Slide: Checkpoints detail
% -----------------------------------------------------------------------------
\begin{frame}[fragile]{Checkpoints: Bounding Recovery Time}
\begin{lstlisting}[style=ascii]
  WITHOUT checkpoints            WITH checkpoints
  -----------------------        -----------------------
  WAL start       crash          CKPT_1      CKPT_2  crash
     |               |              |           |       |
     v               v              v           v       v
-----+~~~~~~~~~~~~~~~+-->      -----+~~~~~~~~~~~+~~~~~~~+-->
     |                              |           |       |
     |<-- replay everything ------->|           |       |
     |  (unbounded recovery time)               |       |
                                                |<----->|
                                                 replay only
                                                 this window

  At each checkpoint:
  (1) flush all dirty pages from shared buffers to data files
  (2) write CKPT record to WAL with redo_lsn

  On crash: replay only WAL after the last redo_lsn
\end{lstlisting}

{\small Triggered by: \texttt{checkpoint\_timeout} (default 5\,min)
or \texttt{max\_wal\_size} exceeded.\quad
Larger \texttt{max\_wal\_size} = better throughput, longer recovery.}
\end{frame}

% -----------------------------------------------------------------------------
% Slide: Exercise 1.3 cue
% -----------------------------------------------------------------------------
\begin{frame}{Exercise 1.3 --- Checkpoints \& Write Path}

  \begin{block}{}
    \faTerminal\enspace\texttt{modules/01-wal-durability/exercises.sql}\\[0.5em]
    Run an explicit \texttt{CHECKPOINT}, insert rows, watch the recovery
    horizon grow in \texttt{pg\_stat\_wal}.\\[0.3em]
    Checkpoint again; check \texttt{pg\_stat\_bgwriter} for
    \texttt{buffers\_backend} --- the write-pressure signal.
  \end{block}

\end{frame}

% -----------------------------------------------------------------------------
% Slide: Full-Page Writes
% -----------------------------------------------------------------------------
\begin{frame}{Full-Page Writes}

  \begin{block}{}
    WAL records the \textit{change} --- a delta.  Recovery replays
    that delta onto the page on disk.\\[0.4em]
    But what if the page on disk was only half-written when the
    crash happened?
  \end{block}

\end{frame}

% -----------------------------------------------------------------------------
% Slide: Full-Page Writes — Backup Blocks
% -----------------------------------------------------------------------------
\begin{frame}[fragile]{Full-Page Writes: Backup Blocks}

  \begin{columns}[T]
    \column{0.44\textwidth}
    \textbf{The fix: backup blocks (full-page images)}
    \begin{itemize}\small
      \item On the \textit{first write to any page after each
            checkpoint}, PostgreSQL embeds the entire \textbf{8\,kB
            page} in the WAL record
      \item Subsequent writes to the same page before the next
            checkpoint are delta-only
      \item Recovery overwrites the corrupted page wholesale from
            the backup block, then replays later deltas cleanly
    \end{itemize}
    {\small \texttt{full\_page\_writes = on} by default.
    Overhead is visible in \texttt{pg\_stat\_wal.wal\_fpi}.}

    \column{0.54\textwidth}
\begin{lstlisting}[style=ascii]
  --[CKPT]------+------------+--> WAL
                |            |
           1st mod      2nd mod
             of page     of page
                |            |
                v            v
           [FPI: full    [delta:
            8 kB page]    small]

  OS crash while writing dirty page:
  +----------+
  | TABLE_A  |  torn: partial write
  | ????     |  --> corrupted on disk
  +----------+

  Recovery:
    FPI  -> overwrites torn page (safe)
    delta -> replays cleanly on good page
\end{lstlisting}
  \end{columns}
\end{frame}

% -----------------------------------------------------------------------------
% Slide: Full-Page Images motivation
% -----------------------------------------------------------------------------
\begin{frame}{Full-Page Images: The Checkpoint Tax}

  \begin{block}{}
    After each checkpoint, the first write to any dirty page embeds
    the \textit{entire} 8\,kB page in the WAL record --- not just the delta.
    This guarantees recovery even from a torn page.\\[0.4em]
    There are no free lunches.  What does this cost?
  \end{block}

\end{frame}

% -----------------------------------------------------------------------------
% Slide: Full-Page Images — cost in practice (ASCII + observe + amortization)
% -----------------------------------------------------------------------------
\begin{frame}[fragile]{Full-Page Images: The Cost in Practice}

  \begin{columns}[T]
    \column{0.44\textwidth}
    \begin{itemize}\small
      \item \textbf{Observe:} \texttt{wal\_fpi} in \texttt{pg\_stat\_wal}
            counts FPI records --- it spikes right after each checkpoint;
            that is expected
      \item \textbf{Amortization:} each page pays the FPI cost
            \textit{once} per checkpoint interval; every subsequent
            write to the same page is a small delta until the
            next checkpoint resets the clock
    \end{itemize}

    \column{0.54\textwidth}
\begin{lstlisting}[style=ascii]
-- writes to the same page over time:

CKPT          CKPT
 |             |
 +--w1--w2--w3-+--w1--w2--w3-->
    |   |   |     |   |   |
  8kB 100B 100B 8kB 100B 100B
  FPI  d    d   FPI  d    d

FPI = full 8kB page in WAL record
d   = delta (only the changed bytes)

-- each page pays FPI cost once per
-- checkpoint interval, not per write
\end{lstlisting}
  \end{columns}

\end{frame}

% -----------------------------------------------------------------------------
% Slide: Full-Page Images — managing the overhead (production tips)
% -----------------------------------------------------------------------------
\begin{frame}{Full-Page Images: Managing the Overhead}

  \begin{itemize}
    \item \textbf{\texttt{wal\_compression = lz4}} (or \texttt{zstd})
          compresses the FPI payload in-place --- meaningfully reduces WAL
          volume on compressible pages with little CPU overhead; safe to
          enable with no schema changes
    \item \textit{Never} disable \texttt{full\_page\_writes} unless your
          filesystem or storage layer guarantees atomic 8\,kB writes ---
          almost none do in practice, and a torn page with only delta
          records in WAL is unrecoverable
  \end{itemize}

\end{frame}

% -----------------------------------------------------------------------------
% Slide: Exercise 1.1 cue
% -----------------------------------------------------------------------------
\begin{frame}{Exercise 1.1 --- Observing WAL Generation}

  \begin{block}{}
    \faTerminal\enspace\texttt{modules/01-wal-durability/exercises.sql}\\[0.5em]
    Snapshot \texttt{pg\_current\_wal\_lsn()}, insert rows, measure the
    delta.\\[0.3em]
    Force a \texttt{CHECKPOINT}, update rows, watch \texttt{wal\_fpi}
    spike in \texttt{pg\_stat\_wal}.
  \end{block}

\end{frame}

% -----------------------------------------------------------------------------
% Slide: WAL Writer motivation
% -----------------------------------------------------------------------------
\begin{frame}{The WAL Writer Process}

  \begin{block}{}
    Every committed transaction must flush WAL to disk.\\[0.4em]
    With hundreds of concurrent commits per second ---
    how do you avoid turning that into hundreds of sequential
    \texttt{fsync()} calls?
  \end{block}

\end{frame}

% -----------------------------------------------------------------------------
% Slide: WAL Writer detail
% -----------------------------------------------------------------------------
\begin{frame}{The WAL Writer Process: Group Commit}

  \textbf{What it does:}
  \begin{itemize}
    \item Background process; wakes every \texttt{wal\_writer\_delay}
          (default 200\,ms)
    \item Flushes pending WAL from shared buffers to \texttt{pg\_wal/}
    \item Also woken immediately by backends waiting on
          \texttt{synchronous\_commit = on}
  \end{itemize}

  \vspace{0.5em}
  \textbf{Group commit:}
  \begin{itemize}
    \item Multiple concurrent \texttt{COMMIT}s accumulated in the WAL buffer
          are flushed in a single \texttt{write()} + \texttt{fsync()}
    \item With high concurrency: \texttt{wal\_write} $\ll$ number of commits
          --- a real saving on spinning disk, meaningful on NVMe
  \end{itemize}

  \vspace{0.3em}
  {\small Observable in \texttt{pg\_stat\_wal}: compare \texttt{wal\_write}
  and \texttt{wal\_sync} counts against the commit count.}

  \vspace{0.4em}
  \alert{The catch:} each committing backend still \textit{waits} for the
  shared fsync to complete before \texttt{COMMIT} returns --- the cost is
  amortized across concurrent transactions, not eliminated.

\end{frame}

% -----------------------------------------------------------------------------
% Slide: synchronous_commit motivation
% -----------------------------------------------------------------------------
\begin{frame}{Durability vs Performance tradeoff}

  \begin{block}{}
    Durability has a cost: every \texttt{COMMIT} waits for WAL to be
    fsynced to disk before returning to the client.\\[0.4em]
    Can the database offer a tunable knob for trading off RPO\footnote{RPO (Recovery Point Objective) --- the maximum committed data loss a workload can tolerate after a crash.} and
    performance for different kinds of workloads?
  \end{block}

\end{frame}

% -----------------------------------------------------------------------------
% Slide: synchronous_commit detail
% -----------------------------------------------------------------------------
\begin{frame}{\texttt{synchronous\_commit}: The Durability Dial}

  {\small Set per-session or per-transaction. Cluster default: \texttt{on}.}

  \vspace{0.4em}
  \begin{tabularx}{\textwidth}{l X l}
    \toprule
    \textbf{Setting} & \textbf{COMMIT returns when \ldots} & \textbf{Max data loss} \\
    \midrule
    \texttt{on}             & WAL written \textit{and} fsynced locally          & none \\[3pt]
    \texttt{remote\_apply} & WAL applied on standby before COMMIT returns      & none on primary \\[3pt]
    \texttt{remote\_write} & WAL written to standby OS buffer                  & standby hard crash \\[3pt]
    \texttt{local}         & WAL fsynced locally only (standby may lag)        & none on primary \\[3pt]
    \texttt{off}           & WAL appended to shared buffer                     & $\leq$ \texttt{wal\_writer\_delay} \\
    \bottomrule
  \end{tabularx}

  \vspace{0.5em}
  \alert{Key:} \texttt{off} does \textit{not} risk corruption.
  The database always stays \textbf{consistent} --- you only risk losing
  the last $\leq$\,200\,ms of committed data.
\end{frame}

% -----------------------------------------------------------------------------
% Slide: Exercise 1.2 cue
% -----------------------------------------------------------------------------
\begin{frame}{Exercise 1.2 --- \texttt{synchronous\_commit}}

  \begin{block}{}
    \faTerminal\enspace\texttt{modules/01-wal-durability/exercises.sql}\\[0.5em]
    Run the benchmark.  Watch \texttt{wal\_write} and \texttt{wal\_sync}
    in \texttt{pg\_stat\_wal}.\\[0.3em]
    The sync \textit{count} is the signal --- not the wall-clock time.
  \end{block}

\end{frame}

% -----------------------------------------------------------------------------
% Slide: Module 1 Exercises
% -----------------------------------------------------------------------------
\begin{frame}{Module 1: Hands-On Exercises}
  \begin{columns}[T]
    \column{0.48\textwidth}
    \textbf{Exercise 1.1 --- Observing WAL generation}
    \begin{itemize}\small
      \item Snapshot LSN with \texttt{pg\_current\_wal\_lsn()}
      \item Insert rows; measure WAL produced
      \item Force a checkpoint; update existing rows
      \item Watch \texttt{wal\_fpi} spike (FPI tax)
    \end{itemize}

    \vspace{0.4em}
    \textbf{Exercise 1.2 --- \texttt{synchronous\_commit} benchmark}
    \begin{itemize}\small
      \item Run inserts with \texttt{on} vs \texttt{off}
      \item Compare \texttt{wal\_write} and \texttt{wal\_sync} deltas
      \item Understand why the sync count is the real signal
    \end{itemize}

    \column{0.48\textwidth}
    \textbf{Exercise 1.3 --- Checkpoints \& write path}
    \begin{itemize}\small
      \item Establish baseline with explicit \texttt{CHECKPOINT}
      \item Insert rows; watch recovery horizon grow
      \item Checkpoint; watch it shrink
      \item Read \texttt{pg\_stat\_bgwriter} for I/O accounting
      \item Discuss \texttt{buffers\_backend} in production
    \end{itemize}

    \vspace{0.5em}
    {\small\faTerminal\enspace
    \texttt{modules/01-wal-durability/exercises.sql}}
  \end{columns}
\end{frame}

% =============================================================================
\section{Module 2 --- MVCC \& Concurrency}
% =============================================================================

% -----------------------------------------------------------------------------
% Slide: MVCC motivation
% -----------------------------------------------------------------------------
\begin{frame}{MVCC: Multi-Version Concurrency Control}

  \begin{block}{}
    In a lock-based system, a writer holds an exclusive lock on a row
    while it works.  Every concurrent reader must wait.\\[0.4em]
    At high concurrency, that wait queue becomes the bottleneck.
    Can we do better?
  \end{block}

\end{frame}

% -----------------------------------------------------------------------------
% Slide: MVCC — The Core Idea
% -----------------------------------------------------------------------------
\begin{frame}[fragile]{MVCC: The Core Idea}

\begin{lstlisting}[style=ascii]
heap page after UPDATE (txid 200):

  tuple_1: xmin=199  xmax=200  name='Jekyll'  <- dead (updated)
  tuple_2: xmin=200  xmax=0    name='Hyde'    <- live

Both versions sit on the same page until VACUUM reclaims the dead one.
A reader with a snapshot taken before txid 200 sees 'Jekyll'.
A reader after sees 'Hyde'.  Neither blocks the other.
\end{lstlisting}

  \vspace{0.5em}
  \textbf{Readers never block writers.  Writers never block readers.}

  \vspace{0.3em}
  {\small PostgreSQL keeps old versions \textit{in the heap page itself}.
  Simpler recovery --- but dead tuples accumulate until VACUUM.}

\end{frame}

% -----------------------------------------------------------------------------
% Slide: xmin and xmax
% -----------------------------------------------------------------------------
\begin{frame}[fragile]{Tuple Versions: \texttt{xmin} and \texttt{xmax}}

  \begin{columns}[T]
    \column{0.48\textwidth}
    Every heap tuple carries two transaction IDs:
    \begin{itemize}\small
      \item \textbf{\texttt{xmin}} --- txid that \textit{created} this version
      \item \textbf{\texttt{xmax}} --- txid that \textit{deleted or superseded} it
            (\texttt{0} = still live)
      \item \textbf{\texttt{ctid}} --- physical pointer to the newer version
    \end{itemize}

    \vspace{0.4em}
    These are system columns --- always present, normally hidden
    from \texttt{SELECT *}.  Full tuple structure
    (\texttt{t\_infomask}, ItemId layout) is in the \textbf{Appendix}.

    \column{0.50\textwidth}
\begin{lstlisting}[style=ascii]
Before UPDATE:
 xmin | xmax | ctid  | balance
------+------+-------+---------
  200 |    0 | (0,1) | 9000.00

Inside txid 201 after UPDATE:
 xmin | xmax | ctid  | balance
------+------+-------+---------
  200 |  201 | (0,2) | 9000.00  <- old
  201 |    0 | (0,2) | 8000.00  <- new

xmax=201 on old tuple signals
"deleted by txid 201".
ctid chain leads to live version.
\end{lstlisting}
  \end{columns}
\end{frame}

% -----------------------------------------------------------------------------
% Slide: Transaction Snapshots motivation
% -----------------------------------------------------------------------------
\begin{frame}{Transaction Snapshots}

  \begin{block}{}
    Multiple versions of the same row coexist on disk.
    \texttt{xmin} and \texttt{xmax} say who created them.\\[0.4em]
    But how does a transaction decide which version it is
    allowed to see?
  \end{block}

\end{frame}

% -----------------------------------------------------------------------------
% Slide: Transaction Snapshots detail
% -----------------------------------------------------------------------------
\begin{frame}[fragile]{Transaction Snapshots}

  \begin{columns}[T]
    \column{0.50\textwidth}
    A \textbf{snapshot} records which transactions were active at a
    point in time.  Format: \texttt{xmin:xmax:xip\_list}

    \begin{itemize}\small
      \item \textbf{xmin}: all txids $<$ xmin are already committed
      \item \textbf{xmax}: all txids $\geq$ xmax have not yet started
      \item \textbf{xip\_list}: txids in between that are still running
    \end{itemize}

    \vspace{0.4em}
    Any txid in \texttt{xip\_list} is treated as IN\_PROGRESS even if
    it has since committed --- this is what makes REPEATABLE READ work.

    \column{0.48\textwidth}
\begin{lstlisting}[style=ascii]
Snapshot: 100:104:100,102
  xmin=100, xmax=104
  active: txids 100 and 102

READ COMMITTED:
  new snapshot per statement
  --> committed changes become
      visible mid-transaction

REPEATABLE READ:
  snapshot at first statement
  --> same view for entire txn,
      no matter what commits
\end{lstlisting}
  \end{columns}
\end{frame}

% -----------------------------------------------------------------------------
% Slide: Visibility Check Rules
% -----------------------------------------------------------------------------
\begin{frame}[fragile]{Visibility Check Rules}

  {\small For each tuple, PostgreSQL tests \texttt{xmin}/\texttt{xmax}
  against the current snapshot:}

  \vspace{0.3em}
\begin{lstlisting}[style=ascii]
t_xmin = ABORTED?             --> invisible  (rollback)

t_xmin = IN_PROGRESS?
  same txid, t_xmax = 0       --> visible    (I just inserted it)
  otherwise                   --> invisible  (dirty read blocked)

t_xmin = COMMITTED?
  t_xmin active in snapshot   --> invisible  (snapshot fence)
  t_xmax = 0 or ABORTED       --> visible    (live tuple)
  t_xmax IN_PROGRESS != me    --> visible    (concurrent delete, not done)
  t_xmax COMMITTED,
    not active in snapshot     --> invisible  (deleted and gone)
\end{lstlisting}

  \vspace{0.3em}
  {\small The snapshot acts as a \textbf{fence}: transactions that committed
  after your snapshot was taken are treated as still in-progress,
  preserving a consistent view.  Dirty reads are structurally impossible.}

\end{frame}

% -----------------------------------------------------------------------------
% Slide: Isolation Levels
% -----------------------------------------------------------------------------
\begin{frame}[fragile]{Isolation Levels}

  \begin{columns}[T]
    \column{0.48\textwidth}
    \textbf{READ COMMITTED} (default)\\
    {\small Fresh snapshot on each statement.}

\begin{lstlisting}[style=ascii]
Session A (READ COMMITTED):
  T1: SELECT count(*) -> 10
  -- Session B inserts + commits
  T2: SELECT count(*) -> 11
      (new snapshot at T2)
\end{lstlisting}

    \vspace{0.5em}
    \textbf{REPEATABLE READ}\\
    {\small Snapshot fixed at first statement.}

\begin{lstlisting}[style=ascii]
Session A (REPEATABLE READ):
  T1: SELECT count(*) -> 10
  -- Session B inserts + commits
  T2: SELECT count(*) -> 10
      (snapshot unchanged)
\end{lstlisting}

    \column{0.48\textwidth}
    \vspace{0.5em}
    \begin{tabular}{lll}
      \toprule
      \textbf{Level} & \textbf{NR Read} & \textbf{Phantom} \\
      \midrule
      READ COMMITTED  & possible  & possible  \\
      REPEATABLE READ & prevented & prevented \\
      SERIALIZABLE    & prevented & prevented \\
      \bottomrule
    \end{tabular}

    \vspace{0.4em}
    {\small Dirty reads are impossible at all levels in PostgreSQL.\\[0.3em]
    REPEATABLE READ also prevents phantoms (SI property).\\
    SSI (SERIALIZABLE) not covered today.}
  \end{columns}
\end{frame}

% -----------------------------------------------------------------------------
% Slide: Lost Update motivation
% -----------------------------------------------------------------------------
\begin{frame}{The Lost Update Problem}

  \begin{block}{}
    MVCC guarantees readers never block writers and vice versa.\\[0.4em]
    But what about two transactions that both read the same value
    and then update it?  Does MVCC protect against that?
  \end{block}

\end{frame}

% -----------------------------------------------------------------------------
% Slide: Lost Update detail
% -----------------------------------------------------------------------------
\begin{frame}[fragile]{The Lost Update Problem}

\begin{lstlisting}[style=ascii]
Session A                          Session B
---------                          ---------
BEGIN;
SELECT balance --> 5000
                                   BEGIN;
                                   SELECT balance --> 5000
                                   UPDATE SET balance = 4700  -- debit 300
                                   COMMIT;
UPDATE SET balance = 4500  -- debit 500 from STALE read
COMMIT;

Final balance: 4500  (Session B's debit is GONE)
\end{lstlisting}

  \vspace{0.4em}
  \begin{columns}[T]
    \column{0.48\textwidth}
    \textbf{Dangerous:} application read-modify-write\\
    {\small Read value, compute new value in app, write it back.
    The write silently overwrites concurrent updates.}

    \column{0.48\textwidth}
    \textbf{Safe:} relative \texttt{UPDATE}\\
    {\small \texttt{SET balance = balance - 500}\\
    pushes arithmetic into the database --- atomic, no stale read.}
  \end{columns}

\end{frame}

% -----------------------------------------------------------------------------
% Slide: SELECT FOR UPDATE motivation
% -----------------------------------------------------------------------------
\begin{frame}{\texttt{SELECT FOR UPDATE}}

  \begin{block}{}
    A relative \texttt{UPDATE} is safe for simple arithmetic.
    But some workflows must read a value first --- check a limit,
    compute a fee, apply business logic --- before writing.\\[0.4em]
    How do you make that read-then-write sequence atomic under
    concurrent writers?
  \end{block}

\end{frame}

% -----------------------------------------------------------------------------
% Slide: SELECT FOR UPDATE detail
% -----------------------------------------------------------------------------
\begin{frame}[fragile]{\texttt{SELECT FOR UPDATE}}

  \begin{columns}[T]
    \column{0.50\textwidth}
    \texttt{SELECT \ldots\ FOR UPDATE} acquires a row-level
    \textbf{ExclusiveLock} before the application reads the value.
    A concurrent session attempting the same lock \textbf{blocks}
    until the first transaction commits or rolls back.

    \vspace{0.5em}
    \textit{First-updater-win}: at REPEATABLE READ, the waiting
    session aborts with a serialization error rather than silently
    overwriting.  The application must retry.

    \column{0.48\textwidth}
\begin{lstlisting}[style=ascii]
Session A              Session B
---------              ---------
BEGIN;
SELECT ... FOR UPDATE
--> ExclusiveLock held
                       BEGIN;
                       SELECT ... FOR UPDATE
                       --> BLOCKS
UPDATE ...
COMMIT;
--> lock released
                       --> unblocks,
                           sees post-commit value
                       UPDATE ...
                       COMMIT;
\end{lstlisting}
    {\small Inspect with \texttt{pg\_locks} and
    \texttt{pg\_blocking\_pids()} in the exercises.}
  \end{columns}
\end{frame}

% -----------------------------------------------------------------------------
% Slide: Module 2 Exercises
% -----------------------------------------------------------------------------
\begin{frame}{Module 2: Hands-On Exercises}
  \begin{columns}[T]
    \column{0.48\textwidth}
    \textbf{Exercise 2.1 --- Snapshot isolation}
    \begin{itemize}\small
      \item INSERT in Session A, don't commit
      \item Session B cannot see it
      \item Commit; Session B's next SELECT sees it
    \end{itemize}

    \vspace{0.4em}
    \textbf{Exercise 2.2 --- Isolation levels}
    \begin{itemize}\small
      \item Session A in REPEATABLE READ: snapshot fixed
      \item Session B inserts and commits
      \item Session A re-counts: same result
      \item Repeat with READ COMMITTED: different result
    \end{itemize}

    \column{0.48\textwidth}
    \textbf{Exercise 2.3 --- Lost update}
    \begin{itemize}\small
      \item Race two sessions on the same row
      \item Compare relative vs application read-modify-write
      \item Observe which debit survives
    \end{itemize}

    \vspace{0.4em}
    \textbf{Exercise 2.4 --- \texttt{SELECT FOR UPDATE}}
    \begin{itemize}\small
      \item Session A locks a row; Session B blocks
      \item Inspect \texttt{pg\_locks} and \texttt{pg\_blocking\_pids()}
      \item Session A commits; B unblocks on the post-commit value
    \end{itemize}

    \vspace{0.4em}
    {\small\faTerminal\enspace
    \texttt{modules/02-mvcc-concurrency/session\_\{a,b\}.sql}}
  \end{columns}
\end{frame}

% -----------------------------------------------------------------------------
% Slide: VACUUM motivation
% -----------------------------------------------------------------------------
\begin{frame}{VACUUM}

  \begin{block}{}
    Every \texttt{UPDATE} leaves a dead tuple on the heap page.
    Every \texttt{DELETE} leaves one too.\\[0.4em]
    Dead tuples accumulate.  Concurrent readers may still need an old
    version, so the database cannot reclaim it immediately.\\[0.4em]
    Who cleans this up --- and when?
  \end{block}

\end{frame}

% -----------------------------------------------------------------------------
% Slide: VACUUM solution
% -----------------------------------------------------------------------------
\begin{frame}{VACUUM}

  \begin{block}{}
    \textbf{Autovacuum} runs in the background, scanning tables
    periodically and reclaiming space occupied by dead tuples that
    no active transaction can see any more.
  \end{block}

\end{frame}

% -----------------------------------------------------------------------------
% Slide: VACUUM overview
% -----------------------------------------------------------------------------
\begin{frame}[fragile]{VACUUM: Overview}

  \begin{columns}[T]
    \column{0.52\textwidth}
    \textbf{What VACUUM does:}
    \begin{itemize}\small
      \item Marks dead tuple slots as reusable (does \textit{not}
            shrink the file)
      \item Updates the \textbf{visibility map} --- pages with no dead
            tuples are flagged and skipped on the next pass
      \item \textbf{Freezes} old tuples, replacing their \texttt{xmin}
            with a special frozen ID to prevent transaction ID
            wraparound
    \end{itemize}

    \vspace{0.4em}
    \textbf{Key signals:}
    \begin{itemize}\small
      \item \texttt{pg\_stat\_user\_tables.n\_dead\_tup} ---
            dead tuple backlog
      \item \texttt{age(datfrozenxid)} in \texttt{pg\_database} ---
            distance from XID wraparound
    \end{itemize}

    \column{0.46\textwidth}
\begin{lstlisting}[style=ascii]
After many UPDATEs:

heap page
+---------------------------+
| dead | dead | live | live |
|  tup |  tup |  tup |  tup |
+---------------------------+
  ^      ^
  VACUUM reclaims these slots;
  space is reused for new tuples

Without VACUUM:
  page fills with dead tuples
  --> new rows spill to new pages
  --> table grows (bloat)
  --> sequential scans slow down
\end{lstlisting}
  \end{columns}

\end{frame}

% =============================================================================
\section{Production Numbers}
% =============================================================================

\begin{frame}{\texttt{synchronous\_commit} on Real Hardware}

  {\small Azure PostgreSQL Flexible Server \textbullet\
  4\,vCPU, 16\,GB RAM \textbullet\
  500\,GB Premium SSDv2 \textbullet\
  \texttt{bench\_inserts(500 rows)}}

  \vspace{0.4em}
  \begin{tabularx}{\textwidth}{X r r}
    \toprule
    & \textbf{\texttt{on}} & \textbf{\texttt{off}} \\
    \midrule
    Wall time           & 950\,ms & 98\,ms \\
    \texttt{wal\_writes} & 504     & 33     \\
    \texttt{wal\_syncs}  & 504     & 5      \\
    \bottomrule
  \end{tabularx}

  \vspace{0.4em}
  \begin{columns}[T]
    \column{0.48\textwidth}
    \begin{itemize}\small
      \item \textbf{9.7$\times$ faster} with \texttt{off} ---
            on NVMe-class storage
      \item Per-commit fsync dominates latency,
            not I/O bandwidth
    \end{itemize}

    \column{0.48\textwidth}
    \begin{itemize}\small
      \item Group commit: 500 commits $\to$ 33 writes,
            5 syncs
      \item \texttt{wal\_syncs} 504 $\to$ 5
            (${\approx}100{\times}$) ---
            the sync count is the signal
    \end{itemize}
  \end{columns}

\end{frame}

% =============================================================================
\section{Takeaways}
% =============================================================================

\begin{frame}{Key Takeaways}

  \begin{block}{WAL \& Durability}
    \begin{itemize}\small
      \item WAL is why PostgreSQL survives crashes:
            \textit{always write the log before the data page}
      \item \texttt{synchronous\_commit} is the highest-leverage throughput
            knob; measure sync \textit{counts}, not just timing
      \item Checkpoints bound recovery time; tune \texttt{max\_wal\_size}
            to balance throughput vs restart latency
      \item Full-page images are unavoidable; reduce overhead with
            \texttt{wal\_compression}
    \end{itemize}
  \end{block}

  \vspace{0.2em}
  \begin{block}{MVCC \& Concurrency}
    \begin{itemize}\small
      \item MVCC keeps old row versions on the heap; readers never block writers
      \item Isolation level controls snapshot timing ---
            READ COMMITTED re-snapshots per statement; REPEATABLE READ holds one
      \item Lost updates are \textit{not} prevented by MVCC: use relative
            \texttt{UPDATE} or \texttt{SELECT FOR UPDATE}
      \item Dead tuples from every \texttt{UPDATE}/\texttt{DELETE} must be
            reclaimed by VACUUM --- left unchecked, this becomes table bloat
    \end{itemize}
  \end{block}
\end{frame}

% -----------------------------------------------------------------------------
% Slide: Thank You
% -----------------------------------------------------------------------------
\begin{frame}[standout]
  \begin{center}
    {\Large \textbf{Thank You!}}\\[0.8em]

    {\normalsize Questions? Let's dig in.}\\[1.2em]

    \href{https://github.com/amitprabhudesai/oltp-workloads-pg-workshop}{%
      \faGithub\enspace\texttt{amitprabhudesai/oltp-workloads-pg-workshop}}\\[0.5em]

    {\small Exercises, devcontainer, and these slides are all in the repo.}\\[0.8em]

    \faEnvelope\enspace\texttt{prabhudesai.amit@gmail.com}
  \end{center}
\end{frame}

% -----------------------------------------------------------------------------
% Slide: References
% -----------------------------------------------------------------------------
\begin{frame}{References}
  \begin{thebibliography}{9}

    \bibitem{suzuki}
      Hironobu Suzuki,
      \textit{The Internals of PostgreSQL},
      \url{https://www.interdb.jp/pg/}.
      Chapter 5 (Concurrency Control) and Chapter 9 (WAL).
      \copyright\ All Rights Reserved, Hironobu Suzuki.
      Used for noncommercial educational purposes per the author's terms;
      see \url{https://www.interdb.jp/pg/} for licence details.

    \bibitem{pgdocs}
      PostgreSQL Global Development Group,
      \textit{PostgreSQL 16 Documentation},
      \url{https://www.postgresql.org/docs/16/}.
      Particularly: \textit{Reliability and the Write-Ahead Log}
      and \textit{WAL Configuration}.

    \bibitem{sosurvey}
      Stack Overflow,
      \textit{Developer Survey 2025 --- Technology},
      \url{https://survey.stackoverflow.co/2025/technology}.
      PostgreSQL \#1 most-used database, three years running (2023--2025).

  \end{thebibliography}
\end{frame}

% =============================================================================
% BONUS SECTION: Write Amplification — Where WAL Meets MVCC
% =============================================================================
\section{Bonus: Write Amplification}

% -----------------------------------------------------------------------------
% Motivation
% -----------------------------------------------------------------------------
\begin{frame}[fragile]{One UPDATE. How Many WAL Records?}

  \begin{block}{}
    MVCC means every \texttt{UPDATE} is a new heap tuple.
    When the updated column is indexed, the index must change too.\\[0.4em]
    Two updates. Same table. Same row. How much WAL does each generate?
  \end{block}

  \vspace{0.5em}
\begin{lstlisting}[style=ascii]
-- A: balance is not indexed
UPDATE accounts SET balance = balance - 100.00 WHERE id = 42;

-- B: owner is indexed
UPDATE accounts SET owner = 'alice_b' WHERE id = 42;
\end{lstlisting}

  \vspace{0.5em}
  {\small One of these generates 3--4$\times$ more WAL than the other.
  Which one, and why?}

\end{frame}

% -----------------------------------------------------------------------------
% Mechanism: HOT vs non-HOT
% -----------------------------------------------------------------------------
\begin{frame}[fragile]{HOT Updates: Skipping the Index WAL}

\begin{lstlisting}[style=ascii]
-- A: balance not indexed -> HOT update
-- Heap: one HOT_UPDATE record (~150 bytes)
-- Btree: nothing

-- B: owner indexed -> must update the index too
-- Heap: one UPDATE record
-- Btree: DELETE old entry + INSERT new entry (~2 records, ~600 bytes total)
\end{lstlisting}

  \vspace{0.5em}
  \textbf{HOT (Heap Only Tuple)} fires when two conditions both hold:
  \begin{enumerate}\small
    \item No indexed column was changed
    \item The new tuple fits on the \textit{same heap page} as the old one
  \end{enumerate}

  \vspace{0.3em}
  Verify with \texttt{pg\_walinspect}:
  \begin{itemize}\small
    \item \texttt{resource\_manager = 'Btree'} records: 0 for HOT, 2 for indexed update
    \item \texttt{record\_type = 'HOT\_UPDATE'} vs \texttt{'UPDATE'} in the Heap record
  \end{itemize}

\end{frame}

% -----------------------------------------------------------------------------
% Practical lever 1: fillfactor
% -----------------------------------------------------------------------------
\begin{frame}[fragile]{Lever 1: \texttt{fillfactor}}

  HOT condition 2 --- new tuple fits on the same page --- fails when the page
  is full. \texttt{fillfactor} reserves space so updates stay in-place.

  \vspace{0.6em}
\begin{lstlisting}[style=ascii]
-- Default: fillfactor = 100 (pages packed full; HOT degrades as table grows)
-- Better for update-heavy tables:
CREATE TABLE accounts (...) WITH (fillfactor = 70);
-- or on an existing table:
ALTER TABLE accounts SET (fillfactor = 70);
VACUUM accounts;   -- rewrites pages at the new fill target
\end{lstlisting}

  \vspace{0.6em}
  \begin{itemize}
    \item 70--80 is a common starting point for tables with frequent \texttt{UPDATE}s
          on non-indexed columns
    \item Trade-off: larger table footprint on disk, fewer pages per scan
    \item Confirm HOT is firing: \texttt{pg\_stat\_user\_tables.n\_tup\_hot\_upd}
          should be a large fraction of \texttt{n\_tup\_upd}
  \end{itemize}

\end{frame}

% -----------------------------------------------------------------------------
% Practical lever 2: index audit
% -----------------------------------------------------------------------------
\begin{frame}[fragile]{Lever 2: Index Audit}

  Every index on a column that gets \texttt{UPDATE}d forces non-HOT writes
  and Btree WAL --- whether or not the index is ever used by a query.

  \vspace{0.6em}
\begin{lstlisting}[style=ascii]
-- Find indexes with zero scans since last stats reset
SELECT schemaname, relname, indexrelname, idx_scan
FROM   pg_stat_user_indexes
WHERE  idx_scan = 0
ORDER  BY schemaname, relname;
\end{lstlisting}

  \vspace{0.6em}
  \begin{itemize}
    \item \texttt{idx\_scan = 0} is a candidate for removal ---
          but check \texttt{pg\_stat\_reset()} history first
    \item Partial indexes (\texttt{WHERE status = 'pending'}) cover fewer rows
          and are HOT-compatible for rows outside the partial condition
    \item Before dropping: \texttt{pg\_get\_wal\_records\_info()} to confirm
          the index is contributing Btree WAL on your hottest \texttt{UPDATE} path
  \end{itemize}

\end{frame}

% =============================================================================
% APPENDIX
% =============================================================================
\appendix

\begin{frame}[standout]
  \begin{center}
    {\Large \textbf{Appendix}}\\[0.5em]
    {\normalsize PostgreSQL internals --- deeper cuts}
  \end{center}
\end{frame}

% -----------------------------------------------------------------------------
% Appendix A1: Internal layout of WAL segment file
% -----------------------------------------------------------------------------
\begin{frame}[fragile]{A1. Internal Layout of a WAL Segment}

  \begin{columns}[T]
    \column{0.48\textwidth}
    A WAL segment is \textbf{16\,MB} by default, subdivided into
    \textbf{8\,kB pages}:
    \begin{itemize}\small
      \item \textbf{Page 0}: \texttt{XLogLongPageHeaderData} ---
            includes \texttt{xlp\_sysid}, segment size, block size;
            used as a cross-check on open
      \item \textbf{Pages 1+}: \texttt{XLogPageHeaderData} ---
            includes \texttt{xlp\_magic}, \texttt{xlp\_tli},
            \texttt{xlp\_pageaddr}, \texttt{xlp\_rem\_len}
      \item XLOG records written sequentially; a record may
            \textit{span} page boundaries
            (\texttt{xlp\_rem\_len} tracks the continuation length)
    \end{itemize}

    \column{0.50\textwidth}
\begin{lstlisting}[style=ascii]
WAL segment (16 MB = 2048 x 8 kB pages)
+------------------------------+
| Page 0 (8 kB)                |
| XLogLongPageHeaderData       |
|  xlp_sysid  xlp_seg_size     |
|  xlp_xlog_blcksz             |
| +---------------------------+|
| | XLOG record 1             ||
| | XLOG record 2             ||
| | XLOG record 3 (partial)   ||
+------------------------------+
| Page 1 (8 kB)                |
| XLogPageHeaderData           |
|  xlp_rem_len (continuation)  |
| +---------------------------+|
| | ... record 3 continued    ||
| | XLOG record 4             ||
+------------------------------+
| Page 2  ...  Page 2047       |
+------------------------------+
\end{lstlisting}
  \end{columns}
\end{frame}

% -----------------------------------------------------------------------------
% Appendix A2: Internal layout of XLOG record
% -----------------------------------------------------------------------------
\begin{frame}[fragile]{A2. Internal Layout of an XLOG Record}

  \begin{columns}[T]
    \column{0.46\textwidth}
    Since PostgreSQL 9.5, XLOG records use a common structured format.

    \vspace{0.3em}
    \textbf{Header} (\texttt{XLogRecord}):
    \begin{itemize}\small
      \item \texttt{xl\_tot\_len} --- total record length
      \item \texttt{xl\_xid} --- transaction ID
      \item \texttt{xl\_prev} --- LSN of previous record
      \item \texttt{xl\_rmid} / \texttt{xl\_info} --- resource manager
            and subtype (e.g.\ \texttt{RM\_HEAP} +
            \texttt{XLOG\_HEAP\_INSERT})
      \item \texttt{xl\_crc} --- integrity check
    \end{itemize}

    \vspace{0.2em}
    {\small Recovery dispatches to the right replay function
    (\texttt{heap\_xlog\_insert}, \texttt{heap\_xlog\_update}, \ldots)
    via \texttt{xl\_rmid} + \texttt{xl\_info}.}

    \column{0.52\textwidth}
\begin{lstlisting}[style=ascii]
Non-backup block (INSERT delta):
+-----------------------+
| XLogRecord (header)   |
|  xl_rmid=RM_HEAP      |
|  xl_info=HEAP_INSERT  |
+-----------------------+
| XLogRecordBlockHeader |
|  relfilenode, fork,   |
|  block_no, data_len   |
+-----------------------+
| XLogRecordDataHeader  |
+-----------------------+
| block data (delta)    |
+-----------------------+
| xl_heap_insert (main) |
+-----------------------+

Backup block: block data replaced by
full 8 kB page (FPI); adds
XLogRecordBlockImageHeader with
length, hole_offset, bimg_info.
\end{lstlisting}
  \end{columns}
\end{frame}

% -----------------------------------------------------------------------------
% Appendix A3: Tuple Structure
% -----------------------------------------------------------------------------
\begin{frame}[fragile,label=frA3]{A3. Tuple Structure}

  \begin{columns}[T]
    \column{0.5\textwidth}
    Each row version on a heap page is a \textbf{HeapTuple}:
    \begin{itemize}\small
      \item \textbf{\texttt{t\_xmin}} --- inserting transaction ID
      \item \textbf{\texttt{t\_xmax}} --- deleting/updating transaction ID (0 = live)
      \item \textbf{\texttt{t\_ctid}} --- pointer to the newer version of this row
            (self-referential if this is the latest)
      \item \textbf{\texttt{t\_infomask}} --- hint bits for commit status
            (avoids repeated clog lookups)
      \item \textbf{\texttt{t\_hoff}} --- offset to user data
    \end{itemize}

    \column{0.48\textwidth}
\begin{lstlisting}[style=ascii]
heap page (8 kB)
+-------------------------------+
| PageHeader                    |
+------+------+------+----------+
| lp1  | lp2  | lp3  | ...      | <- ItemId array
+------+------+------+----------+
|   free space (grows inward)   |
+------+---+---+----------------+
| tuple3 | tuple2 | tuple1      |
+--------+--------+-------------+

Each ItemId:
  lp_off   -- byte offset to tuple
  lp_flags -- Normal/Redirect/Dead/Unused
\end{lstlisting}
  \end{columns}
  {\small \texttt{t\_ctid} forms the version chain: old tuple points to new.
  VACUUM marks dead ItemIds as \texttt{Unused} and reclaims their space.}

  \vspace{0.3em}
  \hyperlink{frTupleInsert}{\beamerreturnbutton{Back}}
\end{frame}

% -----------------------------------------------------------------------------
% Appendix A4: Buffer Manager Structure
% -----------------------------------------------------------------------------
\begin{frame}[fragile]{A4. Buffer Manager Structure}

  \begin{columns}[T]
    \column{0.62\textwidth}
    Three layers between a backend and a page on disk:

\begin{lstlisting}[style=ascii,basicstyle=\ttfamily\tiny]
Buffer Table       Descriptors        Buffer Pool
(tag->buf_id)      (one per slot)     (8 kB slots)

[rel,blk=7]  -->  buf_id=0            slot 0 +--------+
                  tag=[rel,7]                | page 7  |
                  dirty=Y  usage=3           | 8 kB    |
                  refcnt=0                   +--------+

[rel,blk=12] -->  buf_id=1            slot 1 +--------+
                  tag=[rel,12]               | page 12 |
                  dirty=N  usage=1           | 8 kB    |
                  refcnt=2 (pinned)          +--------+

freelist     -->  buf_id=2 (empty)    slot 2 +--------+
                                             | [free]  |
                                             +--------+
\end{lstlisting}

    \column{0.35\textwidth}
    \textbf{Key descriptor fields:}
    \begin{itemize}\small
      \item \textbf{tag} --- \texttt{(relfilenode,}
            \texttt{fork, block\_number)}
      \item \textbf{dirty} --- modified; must reach
            disk before slot is reused
      \item \textbf{refcount} --- pin count;
            \texttt{refcnt\,>\,0} blocks eviction
      \item \textbf{usage\_count} --- access
            frequency; drives \textbf{clock-sweep}
      \item \textbf{freelist} --- empty descriptors
            chained at startup
    \end{itemize}

    \vspace{0.3em}
    {\small \texttt{shared\_buffers}: default 128\,MB;
    typically 25\,\% of RAM in production.}
  \end{columns}
\end{frame}

% -----------------------------------------------------------------------------
% Appendix A5: INSERT / DELETE / UPDATE on the Heap
% -----------------------------------------------------------------------------
\begin{frame}[fragile]{A5. INSERT, DELETE, UPDATE on the Heap}

\begin{lstlisting}[style=ascii]
INSERT:
  tuple_1: xmin=100 xmax=0   ctid=(0,1)  name='Jekyll'
  -- one new tuple, no dead tuple

DELETE (by txid 200):
  tuple_1: xmin=100 xmax=200 ctid=(0,1)  name='Jekyll'
  -- tuple_1 is now "dead" once txid 200 commits
  -- still physically on disk until VACUUM

UPDATE (by txid 201):
  tuple_1: xmin=100 xmax=201 ctid=(0,2)  name='Jekyll'  <- dead
  tuple_2: xmin=201 xmax=0   ctid=(0,2)  name='Hyde'    <- live
  -- UPDATE = logical DELETE + INSERT
  -- Every UPDATE creates one dead tuple
  -- Hot heavy-update tables bloat without VACUUM
\end{lstlisting}

  {\small \texttt{t\_ctid} in tuple\_1 points to tuple\_2, so index scans can
  follow the chain to the live version without re-scanning.
  This is the basis of \textbf{HOT} (Heap Only Tuple) updates.}
\end{frame}

% -----------------------------------------------------------------------------
% Appendix A6: exec_simple_query() — the WAL write path
% -----------------------------------------------------------------------------
\begin{frame}[fragile]{A6. \texttt{exec\_simple\_query()}: WAL Write Path}

\begin{lstlisting}[style=ascii]
exec_simple_query()                         [postgres.c]
  |
  +-- (1) ExtendCLOG()                      [clog.c]
  |        mark txid IN_PROGRESS in CLOG
  |
  +-- (2) heap_insert()                     [heapam.c]
  |        insert tuple into shared buffer page
  |   |
  |   +-- (3) XLogInsert()                  [xloginsert.c]
  |            write Heap/INSERT record to WAL buffer
  |            update page pd_lsn := LSN_1
  |
  +-- (4) finish_xact_command()             [postgres.c]
      |    commit action
      |
      +-- (4a) XLogInsert()                 [xloginsert.c]
      |         write Commit record to WAL buffer (LSN_2)
      |
      +-- (5) XLogWrite()                   [xlog.c]
      |        flush WAL buffer -> WAL segment file
      |        (fsync / fdatasync / O_DSYNC per wal_sync_method)
      |
      +-- (6) TransactionIdCommitTree()     [transam.c]
               mark txid COMMITTED in CLOG
\end{lstlisting}

\end{frame}

% -----------------------------------------------------------------------------
% Appendix A7: WAL as a change stream — logical decoding (concepts)
% -----------------------------------------------------------------------------
\begin{frame}{A7. WAL as a Change Stream: Logical Decoding}

  \textbf{How it works}
  \begin{itemize}
    \item \texttt{wal\_level = logical} records enough context to reconstruct row-level changes
    \item A \textbf{replication slot} pins WAL until the consumer ACKs each LSN — guarantees no gaps
    \item Output plugins decode raw XLOG records into a change stream:\\
          \texttt{test\_decoding} (built-in), \texttt{wal2json}, \texttt{pgoutput} (logical replication)
  \end{itemize}

  \vspace{0.8em}
  \textbf{Production uses}
  \begin{itemize}
    \item CDC pipelines — Debezium, Airbyte, Fivetran
    \item Audit logs without triggers
    \item Live zero-downtime migrations
    \item Event sourcing / outbox pattern
  \end{itemize}

  \vspace{0.8em}
  {\small \textbf{Operational gotcha:} an unconsumed slot holds WAL indefinitely.
  Monitor \texttt{pg\_replication\_slots.wal\_status} and set
  \texttt{max\_slot\_wal\_keep\_size} to bound disk exposure.}

\end{frame}

% -----------------------------------------------------------------------------
% Appendix A7b: Logical decoding — psql output
% -----------------------------------------------------------------------------
\begin{frame}[fragile]{A7b. Logical Decoding: Reading the Change Stream}

\begin{lstlisting}[style=ascii]
-- One-time setup (superuser)
SELECT pg_create_logical_replication_slot('audit', 'test_decoding');

-- peek: non-destructive; get: advances the slot
SELECT lsn, xid, data
FROM pg_logical_slot_peek_changes('audit', NULL, NULL);

 lsn       | xid  | data
-----------+------+----------------------------------------------------------
 0/9368D48 | 1001 | BEGIN 1001
 0/9368DB0 | 1001 | table rootconf.accounts: INSERT: id[bigint]:264
           |      |   owner[text]:'alice'  balance[numeric]:100.00
 0/936BB80 | 1001 | COMMIT 1001
 0/936BB80 | 1002 | BEGIN 1002
 0/936BB80 | 1002 | table rootconf.accounts: UPDATE: id[bigint]:264
           |      |   balance[numeric]:200.00
 0/936BC18 | 1002 | COMMIT 1002
\end{lstlisting}

  \vspace{0.4em}
  {\small Each decoded event carries the LSN, transaction boundary, and full
  column values — exactly what a CDC consumer ingests.}

\end{frame}

% -----------------------------------------------------------------------------
% Appendix A8: Commit Log (clog)
% -----------------------------------------------------------------------------
\begin{frame}[fragile]{A8. Commit Log (clog)}

  \begin{columns}[T]
    \column{0.54\textwidth}
    \textbf{What is clog?}
    \begin{itemize}\small
      \item Stores the \textit{commit status} of every transaction: \\
            IN\_PROGRESS / COMMITTED / ABORTED / SUB\_COMMITTED
      \item Physical location: \texttt{\$PGDATA/pg\_xact/}
      \item 2 bits per transaction; 256 kB file covers 1M transactions
      \item Cached in shared memory (\texttt{pg\_xact} buffer pool)
      \item Visibility checks consult clog to evaluate
            \texttt{t\_xmin}/\texttt{t\_xmax} status
    \end{itemize}

    \column{0.44\textwidth}
\begin{lstlisting}[style=ascii]
pg_xact/0000  (first file)
  txid 0: IN_PROGRESS
  txid 1: COMMITTED
  txid 2: ABORTED
  ...

Each 2-bit entry:
  00 = IN_PROGRESS
  01 = COMMITTED
  10 = ABORTED
  11 = SUB_COMMITTED
\end{lstlisting}
    {\scriptsize Old clog files removed by\\
    VACUUM once all transactions\\
    referencing them are frozen.}
  \end{columns}

  \vspace{0.2em}
  {\small \textbf{Hint bits} (\texttt{t\_infomask}) cache the result of clog
  lookups directly in the tuple header, so subsequent reads skip the clog
  entirely for already-settled transactions.}
\end{frame}

% -----------------------------------------------------------------------------
% Appendix A9: VACUUM — Dead Tuples, Bloat, Wraparound
% -----------------------------------------------------------------------------
\begin{frame}[fragile]{A9. VACUUM: Dead Tuples, Bloat, and Wraparound}

  \begin{columns}[T]
    \column{0.52\textwidth}
    \textbf{Why VACUUM exists:}
    \begin{itemize}\small
      \item Every \texttt{UPDATE}/\texttt{DELETE} leaves a \textit{dead tuple}
            on disk --- MVCC cannot reclaim it immediately
      \item Concurrent readers might still need the old version
      \item \textbf{Autovacuum} (default: on) scans tables periodically and
            reclaims dead tuple space; \texttt{pg\_stat\_user\_tables.n\_dead\_tup}
            shows the backlog
      \item Unchecked: pages fill with dead tuples, table grows
            (\textit{bloat}), sequential scans slow down
    \end{itemize}

    \column{0.46\textwidth}
    \textbf{Transaction ID wraparound:}
    \begin{itemize}\small
      \item txid is a 32-bit counter: wraps at $\sim$4 billion
      \item PostgreSQL uses modular arithmetic with a $2^{31}$ horizon:
            tuples with \texttt{xmin} older than $2^{31}$ become \textit{in the future} ---
            a catastrophic visibility inversion
      \item Fix: VACUUM \textit{freezes} old tuples,
            replacing their \texttt{xmin} with a special frozen txid
            that is always visible
    \end{itemize}
\begin{lstlisting}[style=sql]
-- How far are we from wraparound?
SELECT datname,
  age(datfrozenxid) AS xid_age
FROM pg_database
ORDER BY xid_age DESC;
\end{lstlisting}
  \end{columns}
  {\small Autovacuum warns at 200\,M transactions from wraparound;
  PostgreSQL forces a shutdown at 3\,M to prevent data loss.}
\end{frame}

% =============================================================================
\end{document}
